% --- Archon physics formalization source begin ---
% archon:physics
% archon:covers PhyXMiniProblems/problem_phyx_mini_0847.lean
% archon:source-report reports/phyx_mini/problem_phyx_mini_0847.source.json

\chapter{Physics problem phyx\_mini\_0847}
\label{ch:PhyXMiniProblems_problem_phyx_mini_0847}

\paragraph{Problem source.}
\#\# Question

What is the net electric flux through the torus?

\#\# Answer choices

- A: A: 120N m\textasciicircum{}2/C

- B: B: -100N m\textasciicircum{}2/C

- C: C: 200N m\textasciicircum{}2/C

- D: D: -110N m\textasciicircum{}2/C

\#\# Figure caption (auxiliary; use the image as primary evidence)

The image is a diagram of a torus, which appears as a doughnut-shaped object. Inside the torus, there is a red circle with a plus sign, indicating a positive charge labeled "+100 nC." Outside the torus, a blue circle with a minus sign represents a negative charge labeled "−1 nC (inside)." The illustration uses arrows to connect the charges to their respective labels, indicating their locations relative to the torus.

\#\# Dataset metadata

Electrostatics; Physical Model Grounding Reasoning

\paragraph{Recorded answer/context.}
D: D: -110N m\textasciicircum{}2/C

\paragraph{Figure/image path.}
/root/proposal\_for\_physic/science-mango/phyx\_mini\_run/phyx\_data/test\_image/847.png

\paragraph{Formalization target.}
create a compiling Lean file with sorry bodies at `PhyXMiniProblems/problem\_phyx\_mini\_0847.lean`. The Lean declarations must preserve the physical quantities, dimensions or dimensional roles, figure labels, governing-law hypotheses, and final relation expressed by this problem.
Use Mathlib/Physlib names found through LeanExplore where available. If a physics API is missing, introduce faithful local abstractions rather than scalar placeholder aliases.

\begin{theorem}[Physics formalization target]
\label{thm:physics:phyx_mini_0847:target}
\lean{PhyXMiniProblems.ProblemPhyXMini0847.problem_phyx_mini_0847}
\uses{def:physics:phyx-mini-0847:phyxminiproblems-problemphyxmini0847-electricfluxinnewtonmeterssquaredpercoulomb, def:physics:phyx-mini-0847:phyxminiproblems-problemphyxmini0847-toruselectrostaticssetup, def:physics:phyx-mini-0847:phyxminiproblems-problemphyxmini0847-hasphysicaltorusparameters, def:physics:phyx-mini-0847:phyxminiproblems-problemphyxmini0847-matchessuppliedtorusfigure, def:physics:phyx-mini-0847:phyxminiproblems-problemphyxmini0847-usestextbookvacuumpermittivity, def:physics:phyx-mini-0847:phyxminiproblems-problemphyxmini0847-satisfiestorusgausslaw, def:physics:phyx-mini-0847:phyxminiproblems-problemphyxmini0847-displayedfluxvalue, def:physics:phyx-mini-0847:phyxminiproblems-problemphyxmini0847-isuniqueclosestanswerchoice}
The assigned autoformalize agent should translate this physics problem into Lean declarations in the covered file, with theorem and lemma proof bodies written as `by sorry`.
\end{theorem}
\begin{proof}
This is an autoformalization task, not a proof task. Produce statements that can later be proved by the physics prover without weakening the physical model.
\end{proof}

% --- Archon declaration topology begin ---
\section{Declaration topology}
\begin{definition}[Charge Quantity]
  \label{def:physics:phyx-mini-0847:phyxminiproblems-problemphyxmini0847-chargequantity}
  \lean{PhyXMiniProblems.ProblemPhyXMini0847.ChargeQuantity}
  A signed, unit-independent electric charge
\end{definition}
\begin{proof}
  This is the declared model definition of Charge Quantity.
\end{proof}
\begin{definition}[Length Quantity]
  \label{def:physics:phyx-mini-0847:phyxminiproblems-problemphyxmini0847-lengthquantity}
  \lean{PhyXMiniProblems.ProblemPhyXMini0847.LengthQuantity}
  A nonnegative, unit-independent physical length
\end{definition}
\begin{proof}
  This is the declared model definition of Length Quantity.
\end{proof}
\begin{definition}[vacuum Permittivity Dimension]
  \label{def:physics:phyx-mini-0847:phyxminiproblems-problemphyxmini0847-vacuumpermittivitydimension}
  \lean{PhyXMiniProblems.ProblemPhyXMini0847.vacuumPermittivityDimension}
  Vacuum permittivity has SI dimension charge\textasciicircum{}2 * time\textasciicircum{}2 / (mass * length\textasciicircum{}3)
\end{definition}
\begin{proof}
  This is the declared model definition of vacuum Permittivity Dimension.
\end{proof}
\begin{definition}[Vacuum Permittivity Quantity]
  \label{def:physics:phyx-mini-0847:phyxminiproblems-problemphyxmini0847-vacuumpermittivityquantity}
  \lean{PhyXMiniProblems.ProblemPhyXMini0847.VacuumPermittivityQuantity}
  \uses{def:physics:phyx-mini-0847:phyxminiproblems-problemphyxmini0847-vacuumpermittivitydimension}
  A nonnegative, unit-independent vacuum permittivity
\end{definition}
\begin{proof}
  This is the declared model definition of Vacuum Permittivity Quantity.
\end{proof}
\begin{definition}[electric Flux Dimension]
  \label{def:physics:phyx-mini-0847:phyxminiproblems-problemphyxmini0847-electricfluxdimension}
  \lean{PhyXMiniProblems.ProblemPhyXMini0847.electricFluxDimension}
  Electric flux has SI dimension force * length\textasciicircum{}2 / charge = mass * length\textasciicircum{}3 / (time\textasciicircum{}2 * charge)
\end{definition}
\begin{proof}
  This is the declared model definition of electric Flux Dimension.
\end{proof}
\begin{definition}[Electric Flux Quantity]
  \label{def:physics:phyx-mini-0847:phyxminiproblems-problemphyxmini0847-electricfluxquantity}
  \lean{PhyXMiniProblems.ProblemPhyXMini0847.ElectricFluxQuantity}
  \uses{def:physics:phyx-mini-0847:phyxminiproblems-problemphyxmini0847-electricfluxdimension}
  A signed, unit-independent electric flux
\end{definition}
\begin{proof}
  This is the declared model definition of Electric Flux Quantity.
\end{proof}
\begin{definition}[signed SIReadout]
  \label{def:physics:phyx-mini-0847:phyxminiproblems-problemphyxmini0847-signedsireadout}
  \lean{PhyXMiniProblems.ProblemPhyXMini0847.signedSIReadout}
  Read a signed dimensionful quantity in coherent SI units
\end{definition}
\begin{proof}
  This is the declared model definition of signed SIReadout.
\end{proof}
\begin{definition}[nonnegative SIReadout]
  \label{def:physics:phyx-mini-0847:phyxminiproblems-problemphyxmini0847-nonnegativesireadout}
  \lean{PhyXMiniProblems.ProblemPhyXMini0847.nonnegativeSIReadout}
  Read a nonnegative dimensionful quantity in coherent SI units
\end{definition}
\begin{proof}
  This is the declared model definition of nonnegative SIReadout.
\end{proof}
\begin{definition}[charge In Coulombs]
  \label{def:physics:phyx-mini-0847:phyxminiproblems-problemphyxmini0847-chargeincoulombs}
  \lean{PhyXMiniProblems.ProblemPhyXMini0847.chargeInCoulombs}
  \uses{def:physics:phyx-mini-0847:phyxminiproblems-problemphyxmini0847-chargequantity, def:physics:phyx-mini-0847:phyxminiproblems-problemphyxmini0847-signedsireadout}
  Coulomb readout of a signed electric charge
\end{definition}
\begin{proof}
  This is the declared model definition of charge In Coulombs.
\end{proof}
\begin{definition}[charge In Nanocoulombs]
  \label{def:physics:phyx-mini-0847:phyxminiproblems-problemphyxmini0847-chargeinnanocoulombs}
  \lean{PhyXMiniProblems.ProblemPhyXMini0847.chargeInNanocoulombs}
  \uses{def:physics:phyx-mini-0847:phyxminiproblems-problemphyxmini0847-chargequantity, def:physics:phyx-mini-0847:phyxminiproblems-problemphyxmini0847-chargeincoulombs}
  Nanocoulomb readout used by both charge labels in the figure
\end{definition}
\begin{proof}
  This is the declared model definition of charge In Nanocoulombs.
\end{proof}
\begin{definition}[length In Meters]
  \label{def:physics:phyx-mini-0847:phyxminiproblems-problemphyxmini0847-lengthinmeters}
  \lean{PhyXMiniProblems.ProblemPhyXMini0847.lengthInMeters}
  \uses{def:physics:phyx-mini-0847:phyxminiproblems-problemphyxmini0847-lengthquantity, def:physics:phyx-mini-0847:phyxminiproblems-problemphyxmini0847-nonnegativesireadout}
  Metre readout of a physical length
\end{definition}
\begin{proof}
  This is the declared model definition of length In Meters.
\end{proof}
\begin{definition}[vacuum Permittivity In SI]
  \label{def:physics:phyx-mini-0847:phyxminiproblems-problemphyxmini0847-vacuumpermittivityinsi}
  \lean{PhyXMiniProblems.ProblemPhyXMini0847.vacuumPermittivityInSI}
  \uses{def:physics:phyx-mini-0847:phyxminiproblems-problemphyxmini0847-vacuumpermittivityquantity, def:physics:phyx-mini-0847:phyxminiproblems-problemphyxmini0847-nonnegativesireadout}
  SI readout of vacuum permittivity in C\textasciicircum{}2 / (N m\textasciicircum{}2)
\end{definition}
\begin{proof}
  This is the declared model definition of vacuum Permittivity In SI.
\end{proof}
\begin{definition}[electric Flux In Newton Meters Squared Per Coulomb]
  \label{def:physics:phyx-mini-0847:phyxminiproblems-problemphyxmini0847-electricfluxinnewtonmeterssquaredpercoulomb}
  \lean{PhyXMiniProblems.ProblemPhyXMini0847.electricFluxInNewtonMetersSquaredPerCoulomb}
  \uses{def:physics:phyx-mini-0847:phyxminiproblems-problemphyxmini0847-electricfluxquantity, def:physics:phyx-mini-0847:phyxminiproblems-problemphyxmini0847-signedsireadout}
  Read electric flux in N m\textasciicircum{}2 / C
\end{definition}
\begin{proof}
  This is the declared model definition of electric Flux In Newton Meters Squared Per Coulomb.
\end{proof}
\begin{definition}[Figure Charge]
  \label{def:physics:phyx-mini-0847:phyxminiproblems-problemphyxmini0847-figurecharge}
  \lean{PhyXMiniProblems.ProblemPhyXMini0847.FigureCharge}
  The two point charges distinguished by the labels in the image
\end{definition}
\begin{proof}
  This is the declared model definition of Figure Charge.
\end{proof}
\begin{definition}[Charge Sign Mark]
  \label{def:physics:phyx-mini-0847:phyxminiproblems-problemphyxmini0847-chargesignmark}
  \lean{PhyXMiniProblems.ProblemPhyXMini0847.ChargeSignMark}
  The sign symbol drawn inside a charge marker
\end{definition}
\begin{proof}
  This is the declared model definition of Charge Sign Mark.
\end{proof}
\begin{definition}[Torus Relative Region]
  \label{def:physics:phyx-mini-0847:phyxminiproblems-problemphyxmini0847-torusrelativeregion}
  \lean{PhyXMiniProblems.ProblemPhyXMini0847.TorusRelativeRegion}
  Regions relative to a toroidal Gaussian surface. Only solidInterior is inside the volume bounded by the surface. The central hole belongs to the complement of that volume
\end{definition}
\begin{proof}
  This is the declared model definition of Torus Relative Region.
\end{proof}
\begin{definition}[torus Enclosure Weight]
  \label{def:physics:phyx-mini-0847:phyxminiproblems-problemphyxmini0847-torusenclosureweight}
  \lean{PhyXMiniProblems.ProblemPhyXMini0847.torusEnclosureWeight}
  \uses{def:physics:phyx-mini-0847:phyxminiproblems-problemphyxmini0847-torusrelativeregion}
  The indicator used when summing charge enclosed by a toroidal surface
\end{definition}
\begin{proof}
  This is the declared model definition of torus Enclosure Weight.
\end{proof}
\begin{definition}[Torus Gaussian Surface]
  \label{def:physics:phyx-mini-0847:phyxminiproblems-problemphyxmini0847-torusgaussiansurface}
  \lean{PhyXMiniProblems.ProblemPhyXMini0847.TorusGaussianSurface}
  \uses{def:physics:phyx-mini-0847:phyxminiproblems-problemphyxmini0847-lengthquantity}
  The major and tube radii preserve the toroidal geometric role even though the source supplies no calibrated length scale
\end{definition}
\begin{proof}
  This is the declared model definition of Torus Gaussian Surface.
\end{proof}
\begin{definition}[Torus Charge Figure]
  \label{def:physics:phyx-mini-0847:phyxminiproblems-problemphyxmini0847-toruschargefigure}
  \lean{PhyXMiniProblems.ProblemPhyXMini0847.TorusChargeFigure}
  \uses{def:physics:phyx-mini-0847:phyxminiproblems-problemphyxmini0847-figurecharge, def:physics:phyx-mini-0847:phyxminiproblems-problemphyxmini0847-chargesignmark, def:physics:phyx-mini-0847:phyxminiproblems-problemphyxmini0847-torusrelativeregion}
  Qualitative and numerical data read from the primary image
\end{definition}
\begin{proof}
  This is the declared model definition of Torus Charge Figure.
\end{proof}
\begin{definition}[Torus Electrostatics Setup]
  \label{def:physics:phyx-mini-0847:phyxminiproblems-problemphyxmini0847-toruselectrostaticssetup}
  \lean{PhyXMiniProblems.ProblemPhyXMini0847.TorusElectrostaticsSetup}
  \uses{def:physics:phyx-mini-0847:phyxminiproblems-problemphyxmini0847-chargequantity, def:physics:phyx-mini-0847:phyxminiproblems-problemphyxmini0847-vacuumpermittivityquantity, def:physics:phyx-mini-0847:phyxminiproblems-problemphyxmini0847-electricfluxquantity, def:physics:phyx-mini-0847:phyxminiproblems-problemphyxmini0847-figurecharge, def:physics:phyx-mini-0847:phyxminiproblems-problemphyxmini0847-torusgaussiansurface, def:physics:phyx-mini-0847:phyxminiproblems-problemphyxmini0847-toruschargefigure}
  The independent physical quantities in the electrostatic setup. In particular, netEnclosedCharge and netElectricFlux remain independent until the governing laws below relate them to the individual charges
\end{definition}
\begin{proof}
  This is the declared model definition of Torus Electrostatics Setup.
\end{proof}
\begin{definition}[Has Physical Torus Parameters]
  \label{def:physics:phyx-mini-0847:phyxminiproblems-problemphyxmini0847-hasphysicaltorusparameters}
  \lean{PhyXMiniProblems.ProblemPhyXMini0847.HasPhysicalTorusParameters}
  \uses{def:physics:phyx-mini-0847:phyxminiproblems-problemphyxmini0847-lengthinmeters, def:physics:phyx-mini-0847:phyxminiproblems-problemphyxmini0847-vacuumpermittivityinsi, def:physics:phyx-mini-0847:phyxminiproblems-problemphyxmini0847-toruselectrostaticssetup}
  Positivity and non-self-intersection conditions for the toroidal model
\end{definition}
\begin{proof}
  This is the declared model definition of Has Physical Torus Parameters.
\end{proof}
\begin{definition}[Matches Supplied Torus Figure]
  \label{def:physics:phyx-mini-0847:phyxminiproblems-problemphyxmini0847-matchessuppliedtorusfigure}
  \lean{PhyXMiniProblems.ProblemPhyXMini0847.MatchesSuppliedTorusFigure}
  \uses{def:physics:phyx-mini-0847:phyxminiproblems-problemphyxmini0847-chargeinnanocoulombs, def:physics:phyx-mini-0847:phyxminiproblems-problemphyxmini0847-toruselectrostaticssetup}
  Facts read from the supplied bitmap. The +100 nC charge is in the central hole, not in the solid-torus interior; this topological distinction is the essential figure-derived datum
\end{definition}
\begin{proof}
  This is the declared model definition of Matches Supplied Torus Figure.
\end{proof}
\begin{definition}[Uses Textbook Vacuum Permittivity]
  \label{def:physics:phyx-mini-0847:phyxminiproblems-problemphyxmini0847-usestextbookvacuumpermittivity}
  \lean{PhyXMiniProblems.ProblemPhyXMini0847.UsesTextbookVacuumPermittivity}
  \uses{def:physics:phyx-mini-0847:phyxminiproblems-problemphyxmini0847-vacuumpermittivityinsi, def:physics:phyx-mini-0847:phyxminiproblems-problemphyxmini0847-toruselectrostaticssetup}
  The textbook vacuum-permittivity calibration ε₀ = 8.85 × 10⁻¹² C²/(N m²) used for the numerical estimate
\end{definition}
\begin{proof}
  This is the declared model definition of Uses Textbook Vacuum Permittivity.
\end{proof}
\begin{definition}[Satisfies Torus Gauss Law]
  \label{def:physics:phyx-mini-0847:phyxminiproblems-problemphyxmini0847-satisfiestorusgausslaw}
  \lean{PhyXMiniProblems.ProblemPhyXMini0847.SatisfiesTorusGaussLaw}
  \uses{def:physics:phyx-mini-0847:phyxminiproblems-problemphyxmini0847-chargeincoulombs, def:physics:phyx-mini-0847:phyxminiproblems-problemphyxmini0847-vacuumpermittivityinsi, def:physics:phyx-mini-0847:phyxminiproblems-problemphyxmini0847-electricfluxinnewtonmeterssquaredpercoulomb, def:physics:phyx-mini-0847:phyxminiproblems-problemphyxmini0847-figurecharge, def:physics:phyx-mini-0847:phyxminiproblems-problemphyxmini0847-torusenclosureweight, def:physics:phyx-mini-0847:phyxminiproblems-problemphyxmini0847-toruselectrostaticssetup}
  The governing laws are charge accounting over the bounded solid-torus region and integral Gauss's law Φ\_E = Q\_enclosed / ε₀. They are stated for the general setup and contain neither the requested numerical flux nor an answer choice
\end{definition}
\begin{proof}
  This is the declared model definition of Satisfies Torus Gauss Law.
\end{proof}
\begin{definition}[Answer Choice]
  \label{def:physics:phyx-mini-0847:phyxminiproblems-problemphyxmini0847-answerchoice}
  \lean{PhyXMiniProblems.ProblemPhyXMini0847.AnswerChoice}
  Labels attached to the four displayed flux values
\end{definition}
\begin{proof}
  This is the declared model definition of Answer Choice.
\end{proof}
\begin{definition}[displayed Flux Value]
  \label{def:physics:phyx-mini-0847:phyxminiproblems-problemphyxmini0847-displayedfluxvalue}
  \lean{PhyXMiniProblems.ProblemPhyXMini0847.displayedFluxValue}
  \uses{def:physics:phyx-mini-0847:phyxminiproblems-problemphyxmini0847-answerchoice}
  Numerical value printed for an answer choice, in N m\textasciicircum{}2 / C
\end{definition}
\begin{proof}
  This is the declared model definition of displayed Flux Value.
\end{proof}
\begin{definition}[Is Unique Closest Answer Choice]
  \label{def:physics:phyx-mini-0847:phyxminiproblems-problemphyxmini0847-isuniqueclosestanswerchoice}
  \lean{PhyXMiniProblems.ProblemPhyXMini0847.IsUniqueClosestAnswerChoice}
  \uses{def:physics:phyx-mini-0847:phyxminiproblems-problemphyxmini0847-electricfluxinnewtonmeterssquaredpercoulomb, def:physics:phyx-mini-0847:phyxminiproblems-problemphyxmini0847-toruselectrostaticssetup, def:physics:phyx-mini-0847:phyxminiproblems-problemphyxmini0847-answerchoice, def:physics:phyx-mini-0847:phyxminiproblems-problemphyxmini0847-displayedfluxvalue}
  A choice is uniquely closest to the physical flux when its absolute error is strictly smaller than that of every other displayed value. This definition does not privilege any particular label
\end{definition}
\begin{proof}
  This is the declared model definition of Is Unique Closest Answer Choice.
\end{proof}
\begin{theorem}[enclosed charge of supplied torus figure]
  \label{thm:physics:phyx-mini-0847:phyxminiproblems-problemphyxmini0847-enclosed-charge-of-supplied-torus-figure}
  \lean{PhyXMiniProblems.ProblemPhyXMini0847.enclosed_charge_of_supplied_torus_figure}
  \uses{def:physics:phyx-mini-0847:phyxminiproblems-problemphyxmini0847-chargeinnanocoulombs, def:physics:phyx-mini-0847:phyxminiproblems-problemphyxmini0847-toruselectrostaticssetup, def:physics:phyx-mini-0847:phyxminiproblems-problemphyxmini0847-matchessuppliedtorusfigure, def:physics:phyx-mini-0847:phyxminiproblems-problemphyxmini0847-satisfiestorusgausslaw}
  The supplied figure and charge-accounting law give Q\_enclosed = -1 nC
\end{theorem}
\begin{proof}
  The conclusion follows from the governing relations and figure readouts named in the statement of enclosed charge of supplied torus figure.
\end{proof}
% --- Archon declaration topology end ---
% --- Archon physics formalization source end ---
